\section{Introduction}
\label{sec:intro}

Product formulas are among the most widely used tools in quantum simulation.
The Lie--Trotter formula
\begin{equation}\label{eq:lie-trotter}
  e^{A+B} = \lim_{n\to\infty} \bigl(e^{A/n}\,e^{B/n}\bigr)^n
\end{equation}
decomposes the exponential of a sum into a product of simpler exponentials,
enabling the simulation of composite quantum systems by sequentially
simulating each term of the Hamiltonian~\cite{trotter1959product, lloydUniversalQuantumSimulators1996}.
Higher-order variants---the symmetric Strang splitting~\cite{strang1968construction}
and Suzuki's recursive hierarchy~\cite{suzukiGeneralTheoryFractal1991}---achieve
faster convergence by exploiting symmetry and commutator cancellation.
Rigorous error bounds for these formulas are critical for resource estimation
in quantum algorithms~\cite{childsTheoryTrotterError2021}.

The quality of these error bounds directly affects the efficiency of quantum
algorithms.  The standard bound scales with the product $\norm{A}\norm{B}$,
which treats $A$ and $B$ as unrelated.
Childs et al.~\cite{childsTheoryTrotterError2021} initiated a program of
\emph{commutator-scaling} bounds, showing that the leading error is controlled
by the commutator norm $\norm{\comm{B}{A}}$ rather than the product
$\norm{A}\norm{B}$.
For spatially local Hamiltonians $H = \sum_j h_j$ typical of many-body physics,
most pairs of terms commute, so $\sum_{j<k} \norm{\comm{h_j}{h_k}}$ can be
polynomially smaller than $\sum_{j<k} \norm{h_j}\norm{h_k}$.
These tighter bounds translate directly into fewer Trotter steps---and hence
fewer quantum gates---for a given simulation accuracy.

In this work, we present a complete formalization of the Lie--Trotter product
formula and its higher-order extensions in \Lean, an interactive theorem prover
with a small trusted kernel~\cite{demoura2021lean4}.
The \Mathlib library provides a large corpus of formalized
mathematics~\cite{mathlib2020}, including the normed algebra exponential,
upon which our development builds.

Proofs of product formula error bounds involve long chains of norm estimates
in non-commutative algebras, where each step interacts with the previous through
submultiplicativity, the triangle inequality, and exponential series manipulation.
A single bookkeeping error---a missing factor, a flipped sign in a
commutator---can silently propagate and invalidate the final bound.
Formal verification forces every intermediate step to be explicit and
machine-checked, making it a natural fit for this domain.

Our formalization proves the following hierarchy of results:
\begin{enumerate}
  \item \textbf{First-order Lie--Trotter} (\cref{sec:assembly}):
    $\norm{(e^{A/n}\,e^{B/n})^n - e^{A+B}} \le C/n$
    with explicit constant $C = 2\norm{A}\norm{B}e^{\norm{A}+\norm{B}} + 1$.

  \item \textbf{Strang splitting} (\cref{sec:strang}):
    $\norm{(e^{A/2n}\,e^{B/n}\,e^{A/2n})^n - e^{A+B}} \le C'/n^2$.

  \item \textbf{Multi-operator extensions}:
    $(e^{A_1/n} \cdots e^{A_m/n})^n \to e^{A_1+\cdots+A_m}$
    and the palindromic (Strang) variant, both with explicit error bounds.

  \item \textbf{Commutator-scaling bounds} (\cref{sec:commutator}):
    \begin{align}
      \norm{e^{tB}\,e^{tA} - e^{t(A+B)}}
        &\le \tfrac{1}{2}\norm{\comm{B}{A}}\,t^2\,e^{t(\norm{A}+3\norm{B})},
        \label{eq:comm-first} \\
      \norm{S_2(t) - e^{tH}}
        &\le \bigl(\tfrac{1}{12}\norm{\comm{B}{\comm{B}{A}}}
          + \tfrac{1}{24}\norm{\comm{A}{\comm{A}{B}}}\bigr)\,t^3,
        \label{eq:comm-strang}
    \end{align}
    matching the commutator-scaling program of Childs et al.~\cite{childsTheoryTrotterError2021}.

  \item \textbf{Tighter Strang bound via norm-of-difference}: a strictly
    sharper second-order bound
    $\norm{S_2(t)-e^{tH}} \le \tfrac{\norm{D}}{6}t^3 + \tfrac{T}{4}t^4$
    where $D = \comm{B}{\comm{B}{A/2}} - \comm{A/2}{\comm{A/2}{B}}$ is an
    effective double commutator capturing signed cancellations that the
    triangle-inequality form of \eqref{eq:comm-strang} cannot see.

  \item \textbf{Fourth-order Suzuki $S_4$ with CAS-assisted BCH bounds}
    (\cref{sec:suzuki4}):
    four formalized order-$t^5$ bounds for
    $\norm{S_4(t) - e^{t(A+B)}}$, culminating in a provably
    strictly tighter (9$\times$--64$\times$ per coefficient)
    replacement of Childs's~2021 heuristic prefactors, together with an
    existential-$\delta$ finite-$t$ bound obtained by extending the CAS
    computation to order~$\tau^7$.

  \item \textbf{Fourth-order convergence of $S_4$} (\cref{sec:s4-convergence}):
    compounding the single-step $O(t^5)$ bound over $n$ steps gives
    $\norm{(S_4(t/n))^n - e^{t(A+B)}} \le C/n^4$ for $n$ beyond an explicit
    threshold, and hence $(S_4(t/n))^n \to e^{t(A+B)}$. Together with items~1
    and~2 this completes a verified convergence hierarchy
    $O(1/n) \to O(1/n^2) \to O(1/n^4)$.
\end{enumerate}

The entire development comprises 32 Lean source files totaling over 12\,100 lines,
with zero \texttt{sorry} axioms and zero own theorem-level BCH-interface
axioms on the Lean-Trotter side. The three former BCH-interface axioms
(\texttt{bch\_w4Deriv\_quintic\_level2},
\texttt{bch\_w4Deriv\_level3\_tight},
\texttt{bch\_uniform\_integrated})
are now theorems composing Lean-BCH bridge corollaries with
exp-Lipschitz lifts. As of Lean-BCH revision \texttt{d455ff0}
(2026-05-19), the upstream B1.c quintic axiom has been discharged, so
all $\tau^5$ Suzuki~$S_4$ headlines---together with the standard
Lie--Trotter, Strang, and commutator-scaling theorems---depend only
on the standard foundational axioms \texttt{[propext, Classical.choice,
Quot.sound]}. In particular, axiom inspection on
\texttt{norm\_suzuki4\_childs\_form\_via\_level3} (Level~1 Childs
reproduction), \texttt{norm\_suzuki4\_level3\_bch} (Level~3 tight
$\gamma_i$ prefactors), \texttt{norm\_suzuki4\_level2\_bch} (Level~2
unit-coefficient $\tau^5$ bound), the abstract $S_4$ $O(t^5)$ identity
(\texttt{bch\_iteratedDeriv\_s4Func\_order4}), and \texttt{lie\_trotter}
returns the three foundational axioms only. Two transitive private
axioms remain inside the companion
\href{https://github.com/Jue-Xu/Lean-BCH}{Lean-BCH} project---both
septic stepping stones
(\texttt{symmetric\_bch\_septic\_sub\_poly\_axiom},
\texttt{norm\_septic\_match\_residual\_le\_axiom}). They gate only
the optional Level~4 finite-$t$ uniform refinement with an explicit
$\tau^7$ correction, not the core tight 4th-order bound. All other
Lean-BCH content is fully sorry-free; the three-factor symmetric BCH
cubic identities, the two-factor palindromic quintic remainder, and
the five-factor palindromic quintic identification are theorems
specialized to $\mathbb{K} = \mathbb{R}$ and imported directly.
Our \Lean formalization is available at
\href{https://github.com/Jue-Xu/Lean-Trotter}{\texttt{github.com/Jue-Xu/Lean-Trotter}}.
Our main contributions are:
\begin{itemize}
  \item The \emph{first formalization} of the Lie--Trotter product formula in any
    proof assistant, including its higher-order and multi-operator extensions.

  \item A \emph{Duhamel/FTC-based proof technique} for commutator-scaling bounds
    that avoids ODE uniqueness theory (Gronwall's inequality, Picard iteration),
    requiring only the fundamental theorem of calculus for Bochner integrals.

  \item Machine-checked \emph{commutator-scaling error bounds} at first and
    second order, with double-commutator coefficients matching the
    theoretical literature---and a new tighter-Strang refinement
    via norm-of-difference (\cref{apd:tighter}).

  \item \emph{CAS-assisted BCH bounds for Suzuki $S_4$}: four machine-checked
    bounds---including a rigorously tighter replacement for the
    Childs (2021) prefactors and an existential-$\delta$ finite-$t$ bound
    via an order-$\tau^7$ correction---with all CAS output verified by Lean
    numerical sanity checks.

  \item Identification and filling of several \emph{gaps in \Mathlib}, notably
    the inequality $\norm{e^a} \le e^{\norm{a}}$ for general Banach algebras.
\end{itemize}

To our knowledge, no prior formalization of the Lie--Trotter product formula
exists in any proof assistant.
\Mathlib contains the identity $e^{a+b} = e^a\,e^b$ for commuting elements
(\texttt{exp\_add\_of\_commute}), but the non-commuting convergence
theorem---which is a limit statement, not an algebraic identity---is new.

The remainder of this paper is organized as follows.
\Cref{sec:background} reviews the mathematical background on product formulas
and the Lean/\Mathlib ecosystem.
\Cref{sec:proof} presents the proof architecture, including the modular
decomposition and the Duhamel technique for commutator scaling.
\Cref{sec:formalization} discusses the formalization in \Lean, with
representative code excerpts and design decisions.
\Cref{sec:discussion} reflects on the development process, including
AI-assisted formalization and lessons learned.
\Cref{sec:conclusion} concludes with future directions.
